\documentclass[submit]{harvardml}

\course{CS1810-S26}
\assignment{Homework \#3}
\duedate{March 23, 2026 at 11:59 PM}
\usepackage{float}
\usepackage{common}
\usepackage[OT1]{fontenc}
\usepackage[colorlinks,citecolor=blue,urlcolor=blue]{hyperref}
\usepackage{graphicx}
\usepackage{amsmath}
\usepackage{amssymb}
\usepackage{framed}
\usepackage{color}
\usepackage{listings}
\usepackage{enumitem}
\usepackage{comment}

%%%%%%%%%%%%%%%%%%%%%%%%%%%%%%%%%%%%%%%%%%%
%% Solution environment
\usepackage{xcolor}
\newenvironment{solution}{
    \vspace{2mm}
    \color{blue}\noindent\textbf{Solution}:
}{}
%%%%%%%%%%%%%%%%%%%%%%%%%%%%%%%%%%%%%%%%%%%

\DeclareMathOperator*{\mean}{\mathbb{E}}

\lstset{
  language=Python,
  basicstyle=\ttfamily,
  keywordstyle=\color{blue}\bfseries,
  commentstyle=\color{red},
  stringstyle=\color{green},
  frame=single,
  showstringspaces=false,
}

\definecolor{verbgray}{gray}{0.9}

\lstnewenvironment{csv}{%
  \lstset{backgroundcolor=\color{verbgray},
  frame=single, % hello
  framerule=0pt,
  basicstyle=\ttfamily,
  columns=fullflexible}}{}

\begin{document}

\begin{center}
  {\Large Neural Networks and Kernels}\\
\end{center}

\subsection*{Introduction}

This homework is about neural networks and kernels.

\begin{enumerate}
  \item You'll explore how kernel functions implicitly define feature maps and how kernel ridge regression connects to predictions as weighted sums of training labels.
  \item You'll derive the backpropagation algorithm for a single-hidden-layer neural network for the binary classification task.
  \item You'll discover neural scaling laws empirically and derive optimal compute allocation strategies.
\end{enumerate}

As always, please start early and come to office hours with questions!

\subsection*{Resources and Submission Instructions}
You may want to consider the notes from Section 4 (Richer Features and Neural Networks).

Please type your solutions after the corresponding problems using this \LaTeX\ template, and start each problem on a new page.

Please submit the writeup PDF to the Gradescope assignment `HW3'. Remember to assign pages for each question.  \textbf{You must include any plots in your writeup PDF. }. Please submit your \LaTeX file and code files to the Gradescope assignment `HW3 - Supplemental.' The supplemental files will only be checked in special cases, e.g. honor code issues, etc. Your files should be named in the same way as we provide them in the repository, e.g. \texttt{hw3.pdf}, etc.

\newpage

%%%%%%%%%%%%%%%%%%%%%%%%%%%%%%%%%%%%%%%%%%%%%
% Problem 1
%%%%%%%%%%%%%%%%%%%%%%%%%%%%%%%%%%%%%%%%%%%%%

\begin{problem}[Kernels and Feature Maps, 30 pts]

In this problem, we explore how kernel functions implicitly define
feature maps and how this leads to predictions that are weighted sums of
training labels.

Recall from lecture that a kernel function $K(\boldx, \boldx')$ computes
an inner product in some feature space: $K(\boldx, \boldx') =
\bphi(\boldx)^\trans \bphi(\boldx')$, where $\bphi$ is a feature map.

\begin{enumerate}

  \item[1.] \textbf{Expanding the polynomial kernel.}
  Consider the degree-2 polynomial kernel for $\boldx \in \reals^2$:
  \[
    K(\boldx, \boldx') = (1 + \boldx^\trans \boldx')^2.
  \]

  \begin{enumerate}
    \item Writing $\boldx = (x_1, x_2)$ and $\boldx' = (x_1', x_2')$,
    expand $(1 + x_1 x_1' + x_2 x_2')^2$ and collect terms. Show that the
    result can be written as $\bphi(\boldx)^\trans \bphi(\boldx')$ for
    some feature map $\bphi : \reals^2 \to \reals^6$. State $\bphi$
    explicitly.

    \item Verify your answer with a concrete example: compute both
    $K(\boldx, \boldx')$ and $\bphi(\boldx)^\trans \bphi(\boldx')$ for
    $\boldx = (1, 0)$ and $\boldx' = (1, 1)$.
  \end{enumerate}

  \item[2.] \textbf{Visualizing the feature map.}
  We provide a synthetic dataset of concentric circles: an inner ring
  (class~0) and an outer ring (class~1). See the accompanying notebook
  for the data generation code.

  \begin{enumerate}
    \item Plot the dataset in the original space $\reals^2$.
    Are the data linearly separable? Explain briefly.

    \item Apply the feature map $\bphi$ from Part~1 to the data. Produce
    a 3D scatter plot of $(x_1,\; x_2,\; x_1^2 + x_2^2)$ with points
    colored by class. Are the data linearly separable in feature space?
    Why does this feature map help for concentric circles?

    \item Fit a ridge regression model in the feature space:
    \[
      \boldw = (\bPhi^\trans \bPhi + \lambda \ident)^{-1} \bPhi^\trans \boldy,
    \]
    where $\bPhi \in \reals^{N \times 6}$ is the feature matrix with
    $\bphi(\boldx_n)^\trans$ as its $n$th row and $\boldy \in \{0, 1\}^N$
    are the labels. Use $\lambda = 0.01$. Use pure \texttt{numpy} (no \texttt{sklearn}).
    Plot the decision boundary over the training data.
  \end{enumerate}

  \item[3.] \textbf{From feature space to kernel space.}

  In Part~2c you computed $\boldw$ and made predictions as
  $f(\boldx^*) = \bphi(\boldx^*)^\trans \boldw$, which required
  explicitly constructing the feature matrix $\bPhi$. We now show that
  the same predictions can be written as a weighted sum over training
  points, purely in terms of the kernel function $K$.

  \begin{enumerate}
    \item Let $\boldK \in \reals^{N \times N}$ be the kernel matrix
    with entries $K_{ij} = K(\boldx_i, \boldx_j)$, and define
    the weight vector
    $\balpha = (\boldK + \lambda \ident)^{-1} \boldy$.
    Show that the ridge regression prediction at a new point $\boldx^*$
    can be written as:
    \[
      f(\boldx^*) = \sum_{n=1}^{N} \alpha_n \, K(\boldx_n, \boldx^*).
    \]

    \textit{Hint}: Define $\boldw = \bPhi^\trans \balpha$ and verify
    that this $\boldw$ satisfies the ridge regression normal equations
    $(\bPhi^\trans \bPhi + \lambda \ident)\boldw = \bPhi^\trans \boldy$,
    using $\bPhi \bPhi^\trans = \boldK$.

    \item Implement kernel ridge regression using the polynomial kernel
    $K(\boldx, \boldx') = (1 + \boldx^\trans \boldx')^2$. Compute
    $\balpha = (\boldK + \lambda \ident)^{-1} \boldy$ on the
    training data. Then make a scatter plot of the training data where
    each point is sized proportionally to $|\alpha_n|$.

    \item Which points have the largest $|\alpha_n|$? Where do they tend
    to be located? Briefly explain why.

    \item In kernel ridge regression, every training point contributes to
    the prediction (all $\alpha_n \neq 0$). In a \emph{support vector
    machine} (SVM), only the training points closest to the decision
    boundary have $\alpha_n \neq 0$. These points are called
    \emph{support vectors}.

    Looking at your plot from part~(b), which training points do you
    think would be selected as support vectors? Why might it be
    useful to have most $\alpha_n = 0$?
  \end{enumerate}

  \item[4.] \textbf{The RBF kernel.}

  The radial basis function (RBF) kernel is defined as:
  \[
    K(\boldx, \boldx') = \exp\!\left(-\frac{\|\boldx - \boldx'\|^2}{2\sigma^2}\right).
  \]
  Unlike the polynomial kernel, whose degree limits it to polynomial
  decision boundaries, the RBF kernel can approximate any continuous
  decision boundary. The bandwidth $\sigma$ controls how quickly
  $K(\boldx, \boldx')$ decays with the distance between $\boldx$
  and $\boldx'$.

  \begin{enumerate}
    \item Expand $\|\boldx - \boldx'\|^2$ and use the result to write
    $K(\boldx, \boldx')$ as a product of three exponentials: one
    depending only on $\boldx$, one only on $\boldx'$ and one
    involving $\boldx^\trans \boldx'$.

    \item The other two factors depend on $\boldx$ or $\boldx'$ alone,
    so they act as scalar weights on each input. The factor involving
    $\boldx^\trans \boldx'$ is what couples the two inputs. Expand it
    as a Taylor series.

    \item The kernel trick lets us work with $K$ without ever
    computing $\bphi$ explicitly. Why is this essential for the RBF?

    \item Assume the training points are distinct. With
    $\balpha = (\boldK + \lambda \ident)^{-1} \boldy$ and
    $f(\boldx^*) = \sum_{n=1}^N \alpha_n \, K(\boldx_n, \boldx^*)$,
    show that as $\sigma \to 0$ the model classifies every training
    point correctly but predicts $f(\boldx^*) \to 0$ at any point
    not in the training set.

    \item Show that as $\sigma \to \infty$ the model predicts the
    same value at every point.
  \end{enumerate}

\end{enumerate}
\end{problem}

\newpage

\begin{solution}
\item 1. SubProblem 1

(a)

We expand
\[
K(\boldx,\boldx') = (1 + \boldx^\top \boldx')^2
= (1 + x_1 x_1' + x_2 x_2')^2.
\]

Expanding gives
\[
1 + 2x_1 x_1' + 2x_2 x_2' + x_1^2 x_1'^2 + 2x_1 x_2 x_1' x_2' + x_2^2 x_2'^2.
\]

Define
\[
\bphi(\boldx) =
\begin{bmatrix}
1 \\
\sqrt{2} x_1 \\
\sqrt{2} x_2 \\
x_1^2 \\
\sqrt{2} x_1 x_2 \\
x_2^2
\end{bmatrix}.
\]

Then
\[
\bphi(\boldx)^\top \bphi(\boldx') = K(\boldx,\boldx').
\]

(b)

Let $\boldx = (1,0)$ and $\boldx' = (1,1)$.

\[
K(\boldx,\boldx') = (1 + 1)^2 = 4.
\]

\[
\bphi(1,0) =
\begin{bmatrix}
1 \\ \sqrt{2} \\ 0 \\ 1 \\ 0 \\ 0
\end{bmatrix},
\quad
\bphi(1,1) =
\begin{bmatrix}
1 \\ \sqrt{2} \\ \sqrt{2} \\ 1 \\ \sqrt{2} \\ 1
\end{bmatrix}.
\]

\[
\bphi(\boldx)^\top \bphi(\boldx') = 4.
\]
\item 2. SubProblem 2

(a)

\begin{center}
\includegraphics[width=0.45\textwidth]{output2a.png}
\end{center}

The data consists of two concentric circles, where one class surrounds the other. These data are not linearly separable in $\reals^2$, since any linear boundary would intersect both classes.

(b)

\begin{center}
\includegraphics[width=0.8\textwidth]{output2b.png}
\end{center}

In the feature space $(x_1, x_2, x_1^2 + x_2^2)$, the two classes become separable. The inner circle has smaller values of $x_1^2 + x_2^2$, while the outer circle has larger values, creating separation along the $z$-axis. A linear separator of the form $z = c$ separates the two classes.

(c)

\begin{center}
\includegraphics[width=0.45\textwidth]{output2c.png}
\end{center}

The learned decision boundary is circular and correctly separates the inner and outer circles. This reflects that ridge regression is linear in feature space, but induces a nonlinear (quadratic) boundary in the original space.

\vspace{1em}

\item 3. SubProblem 3


(a)

Let
\[
\balpha = (\boldK + \lambda I)^{-1} \boldy,
\quad
\boldw = \bPhi^\top \balpha.
\]

We verify that $\boldw$ satisfies ridge regression:
\[
(\bPhi^\top \bPhi + \lambda I)\boldw = \bPhi^\top \boldy.
\]

\[
(\bPhi^\top \bPhi + \lambda I)\bPhi^\top \balpha
= \bPhi^\top (\bPhi \bPhi^\top + \lambda I)\balpha
= \bPhi^\top (\boldK + \lambda I)\balpha
= \bPhi^\top \boldy.
\]

Thus $\boldw$ is valid.

Now for a test point $\boldx^*$:
\[
f(\boldx^*) = \bphi(\boldx^*)^\top \boldw
= \bphi(\boldx^*)^\top \bPhi^\top \balpha
= \sum_{n=1}^N \alpha_n K(\boldx_n, \boldx^*).
\]
(b)

We compute
\[
\balpha = (\boldK + \lambda I)^{-1} \boldy,
\]
where
\[
K_{ij} = (1 + \boldx_i^\top \boldx_j)^2.
\]

The scatter plot sizes each point proportionally to $|\alpha_n|$.s

(c)

\begin{center}
\includegraphics[width=0.45\textwidth]{output3c.png}
\end{center}

Points with larger $|\alpha_n|$ are concentrated near the boundary between the two classes. These points contribute most to the prediction function
\[
f(\boldx) = \sum_{n=1}^N \alpha_n K(\boldx_n, \boldx),
\]
since they determine the location of the decision boundary. Points far from the boundary have smaller coefficients and are less influential.

(d)

The support vectors would be the points closest to the decision boundary.

From the plot, these correspond to the inner edge of the outer ring and the outer edge of the inner ring.

These points are most informative because they determine the location of the boundary.

Having most $\alpha_n = 0$ is useful because it makes the model sparse, reducing computation and focusing only on the most important points.

\item 4. SubProblem 4

(a)

\[
\|\boldx - \boldx'\|^2 = \boldx^\top \boldx - 2 \boldx^\top \boldx' + \boldx'^\top \boldx'.
\]

Thus
\[
K(\boldx,\boldx')
=
\exp\left(-\frac{\boldx^\top \boldx}{2\sigma^2}\right)
\exp\left(-\frac{\boldx'^\top \boldx'}{2\sigma^2}\right)
\exp\left(\frac{\boldx^\top \boldx'}{\sigma^2}\right).
\]

(b)

\[
\exp\left(\frac{\boldx^\top \boldx'}{\sigma^2}\right)
=
\sum_{m=0}^{\infty}
\frac{1}{m!}
\left(\frac{\boldx^\top \boldx'}{\sigma^2}\right)^m.
\]

(c)

The feature map for the RBF kernel is infinite-dimensional, so it is not practical to compute explicitly. The kernel trick allows us to compute inner products without constructing the feature map.

(d)

As $\sigma \to 0$:
\[
K(\boldx_i,\boldx_j) \to
\begin{cases}
1 & i=j \\
0 & i \ne j
\end{cases}
\]

Thus, the model memorizes the training data and predicts $0$ elsewhere.

(e)

As $\sigma \to \infty$:
\[
K(\boldx,\boldx') \to 1.
\]

Thus, all predictions become constant.

\end{solution}
\newpage

%%%%%%%%%%%%%%%%%%%%%%%%%%%%%%%%%%%%%%%%%%%%%
% Problem 2
%%%%%%%%%%%%%%%%%%%%%%%%%%%%%%%%%%%%%%%%%%%%%

\begin{problem}[Neural Networks, 20 pts]

In this problem, we will take a closer look at how gradients are calculated for backpropagation with a simple multi-layer perceptron (MLP). The MLP will consist of a first fully connected layer with a sigmoid activation, followed by a one-dimensional, second fully connected layer with a sigmoid activation to get a prediction for a binary classification problem. We use non-linear activation functions as the composition of linear functions is linear. Assume bias has not been merged. Let:
\begin{itemize}
  \item $\bold{W}_1$ be the weights of the first layer, $\bold{b}_1$ be the bias of the first layer.
  \item $\bold{W}_2$ be the weights of the second layer, $\bold{b}_2$ be the bias of the second layer.
\end{itemize}

The described architecture can be written mathematically as: $$\hat{y} = \sigma(\bold{W}_2 \left[\sigma \left(\bold{W}_1 \bold{x} + \bold{b}_1\right)\right] + \bold{b}_2)$$

where $\hat{y}$ is a scalar output of the net when passing in the single datapoint $\bold{x}$ (represented as a column vector), the additions are element wise additions, and the sigmoid is an element wise sigmoid.

\begin{enumerate}
  \item Let:
        \begin{itemize}
          \item $N$ be the number of datapoints we have
          \item $M$ be the dimensionality of the data
          \item $H$ be the size of the hidden dimension of the first layer. Here, hidden dimension is used to describe the dimension of the resulting value after going through the layer. Based on the problem description, the hidden dimension of the second layer should be 1.
        \end{itemize}

        Write out the dimensionality of each of the parameters, and of the intermediate variables:

        \begin{align*}
          \bold{a}_1 & = \bold{W}_1 \bold{x} + \bold{b}_1,
                     & \bold{z}_1 = \sigma(\bold{a}_1)       \\
          a_2        & = \bold{W}_2 \bold{z}_1 + \bold{b}_2,
                     & \hat{y} = z_2 = \sigma(a_2)
        \end{align*}

        and make sure they work with the mathematical operations described above. Examining shapes is one of the key ways to debug your code, and can be done using .shape after any numpy array.

  \item  We will derive the gradients for each of the parameters, which can then be used along with gradient descent to find weights that improve our model's performance. For this question, assume there is only one datapoint $\bold{x}$, and that our loss is $L = -(y \log (\hat{y}) + (1 - y) \log (1 - \hat{y}))$. For all questions, the chain rule will be useful.
        \begin{enumerate}
          \item Find $\frac{\partial L}{\partial b_2}$.

          \item Find $\frac{\partial L}{\partial W_2^h}$, where $W_2^h$ represents the $h$th element of $\bold{W}_2$.

          \item Find $\frac{\partial L}{\partial b_1^h}$, where $b_1^h$ represents the $h$th element of $\bold{b}_1$. (\textbf{Hint}: Note that only the $h$th element of $\bold{a}_1$ and $\bold{z}_1$ depend on $b_1^h$ - this should help you with how to use the chain rule.)

          \item Find $\frac{\partial L}{\partial W_1^{h, j}}$, where  $W_1^{h,j}$ represents the element in row $h$, column $j$ in $\bold{W}_1$.

        \end{enumerate}

\end{enumerate}

\end{problem}

\newpage

\begin{framed}
  \noindent\textbf{Problem 2} (cont.)\\
  \begin{enumerate}
    \setcounter{enumi}{2}

    \item  We now explore an example of forward-mode auto-differentiation. Consider the following
          equation:

          $$
            f(x_1, x_2) = \ln (\sin (x_1)) + x_1 \exp \{ x_2 \}
          $$

          This equation can be split up using intermediate variables $v_1, \dots, v_7$ as follows:

          \begin{align*}
            v_1         & = x_1            \\
            v_2         & = \sin (v_1)     \\
            v_3         & = \ln (v_2)      \\
            v_4         & = x_2            \\
            v_5         & = \exp \{ v_4 \} \\
            v_6         & = v_1v_5         \\
            v_7         & = v_3 + v_6      \\
            f(x_1, x_2) & = v_7
          \end{align*}

          Splitting up the equation like this is very similar to what an auto-differentiation
          library would do. From these equations we can construct a \textit{computational graph}
          where each node of the graph corresponds to an input, an intermediate variable, or
          the output.

          \begin{enumerate}
            \item Let $x_1 = \frac{\pi}{6}$ and $x_2 = 1$. Calculate the values of all the
                  intermediate variables $v_1, \dots v_7$ and $f(x_1,x_2)$.
            \item Calculate the derivative of
                  all of the intermediate variables $v_1, \dots, v_7$ and
                  $f$ with respect to $x_1$ evaluated
                  at $x_1 = \frac{\pi}{6}$ and $x_2 = 1$. Remember to write out the equations before evaluating, e.g.,
                  \[
                    \frac{\partial f(x)}{\partial x} = \frac{\partial f(x)}{\partial g(x)} \frac{\partial g(x)}{\partial x}.
                  \]
          \end{enumerate}
  \end{enumerate}
\end{framed}

\begin{solution}
\begin{enumerate}
    \item[1.] We first determine the shapes of all parameters and intermediate variables.

    Since $\bold{x}$ is a single datapoint with input dimension $M$, we have
    \[
    \bold{x} \in \reals^{M \times 1}.
    \]
    The first layer maps $\reals^M \to \reals^H$, so
    \[
    \bold{W}_1 \in \reals^{H \times M},
    \qquad
    \bold{b}_1 \in \reals^{H \times 1}.
    \]
    Therefore
    \[
    \bold{a}_1 = \bold{W}_1 \bold{x} + \bold{b}_1 \in \reals^{H \times 1},
    \qquad
    \bold{z}_1 = \sigma(\bold{a}_1) \in \reals^{H \times 1}.
    \]

    Since the second layer outputs a scalar, it maps $\reals^H \to \reals$, so
    \[
    \bold{W}_2 \in \reals^{1 \times H},
    \qquad
    \bold{b}_2 \in \reals.
    \]
    Hence
    \[
    a_2 = \bold{W}_2 \bold{z}_1 + \bold{b}_2 \in \reals,
    \qquad
    \hat{y} = z_2 = \sigma(a_2) \in \reals.
    \]

    If we instead process all $N$ datapoints at once, then the input matrix would be in $\reals^{M \times N}$, but for the formulas written here, since we are passing a single column vector $\bold{x}$, the dimensions above are the correct ones.

    \item[2.] We now derive the gradients for one datapoint. The loss is
    \[
    L = -\left(y \log \hat{y} + (1-y)\log(1-\hat{y})\right).
    \]

    First compute the derivative with respect to $a_2$. Using the chain rule,
    \[
    \frac{\partial L}{\partial \hat{y}}
    =
    -\frac{y}{\hat{y}} + \frac{1-y}{1-\hat{y}},
    \qquad
    \frac{\partial \hat{y}}{\partial a_2}
    =
    \hat{y}(1-\hat{y}).
    \]
    Therefore
    \[
    \frac{\partial L}{\partial a_2}
    =
    \frac{\partial L}{\partial \hat{y}}
    \frac{\partial \hat{y}}{\partial a_2}
    =
    \left(-\frac{y}{\hat{y}} + \frac{1-y}{1-\hat{y}}\right)\hat{y}(1-\hat{y})
    =
    \hat{y} - y.
    \]

    \begin{enumerate}
        \item Since
        \[
        a_2 = \bold{W}_2 \bold{z}_1 + b_2,
        \]
        we have
        \[
        \frac{\partial a_2}{\partial b_2} = 1.
        \]
        Thus
        \[
        \frac{\partial L}{\partial b_2}
        =
        \frac{\partial L}{\partial a_2}
        \frac{\partial a_2}{\partial b_2}
        =
        \hat{y} - y.
        \]

        \item Since
        \[
        a_2 = \sum_{h=1}^H W_2^h z_1^h + b_2,
        \]
        we get
        \[
        \frac{\partial a_2}{\partial W_2^h} = z_1^h.
        \]
        Hence
        \[
        \frac{\partial L}{\partial W_2^h}
        =
        \frac{\partial L}{\partial a_2}
        \frac{\partial a_2}{\partial W_2^h}
        =
        (\hat{y} - y) z_1^h.
        \]

        \item We now compute $\frac{\partial L}{\partial b_1^h}$. By the hint, only the $h$th coordinates of $\bold{a}_1$ and $\bold{z}_1$ depend on $b_1^h$. Thus
        \[
        \frac{\partial L}{\partial b_1^h}
        =
        \frac{\partial L}{\partial a_2}
        \frac{\partial a_2}{\partial z_1^h}
        \frac{\partial z_1^h}{\partial a_1^h}
        \frac{\partial a_1^h}{\partial b_1^h}.
        \]
        Now,
        \[
        \frac{\partial L}{\partial a_2} = \hat{y}-y,
        \qquad
        \frac{\partial a_2}{\partial z_1^h} = W_2^h,
        \qquad
        \frac{\partial z_1^h}{\partial a_1^h} = z_1^h(1-z_1^h),
        \qquad
        \frac{\partial a_1^h}{\partial b_1^h} = 1.
        \]
        Therefore
        \[
        \frac{\partial L}{\partial b_1^h}
        =
        (\hat{y}-y)\,W_2^h\,z_1^h(1-z_1^h).
        \]

        \item Finally,
        \[
        a_1^h = \sum_{j=1}^M W_1^{h,j} x^j + b_1^h,
        \]
        so
        \[
        \frac{\partial a_1^h}{\partial W_1^{h,j}} = x^j.
        \]
        Using the same chain rule as before,
        \[
        \frac{\partial L}{\partial W_1^{h,j}}
        =
        \frac{\partial L}{\partial a_2}
        \frac{\partial a_2}{\partial z_1^h}
        \frac{\partial z_1^h}{\partial a_1^h}
        \frac{\partial a_1^h}{\partial W_1^{h,j}}.
        \]
        Substituting each term gives
        \[
        \frac{\partial L}{\partial W_1^{h,j}}
        =
        (\hat{y}-y)\,W_2^h\,z_1^h(1-z_1^h)\,x^j.
        \]
    \end{enumerate}

    \item[3.] 

    The intermediate variables are
    \begin{align*}
        v_1 &= x_1, \\
        v_2 &= \sin(v_1), \\
        v_3 &= \ln(v_2), \\
        v_4 &= x_2, \\
        v_5 &= e^{v_4}, \\
        v_6 &= v_1 v_5, \\
        v_7 &= v_3 + v_6, \\
        f(x_1,x_2) &= v_7.
    \end{align*}

    \begin{enumerate}
        \item Let $x_1 = \frac{\pi}{6}$ and $x_2 = 1$. Then
        \[
        v_1 = \frac{\pi}{6},
        \qquad
        v_2 = \sin\left(\frac{\pi}{6}\right) = \frac{1}{2},
        \qquad
        v_3 = \ln\left(\frac{1}{2}\right) = -\ln 2,
        \]
        \[
        v_4 = 1,
        \qquad
        v_5 = e,
        \qquad
        v_6 = \frac{\pi}{6}e,
        \qquad
        v_7 = -\ln 2 + \frac{\pi e}{6}.
        \]
        Therefore
        \[
        f\left(\frac{\pi}{6},1\right)= -\ln 2 + \frac{\pi e}{6}.
        \]

        \item We now compute all derivatives with respect to $x_1$.

        First,
        \[
        \frac{\partial v_1}{\partial x_1} = 1.
        \]

        Since $v_2 = \sin(v_1)$,
        \[
        \frac{\partial v_2}{\partial x_1}
        =
        \frac{\partial v_2}{\partial v_1}
        \frac{\partial v_1}{\partial x_1}
        =
        \cos(v_1)\cdot 1
        =
        \cos(v_1).
        \]
        Evaluating at $x_1=\frac{\pi}{6}$,
        \[
        \frac{\partial v_2}{\partial x_1}
        =
        \cos\left(\frac{\pi}{6}\right)
        =
        \frac{\sqrt{3}}{2}.
        \]

        Since $v_3 = \ln(v_2)$,
        \[
        \frac{\partial v_3}{\partial x_1}
        =
        \frac{\partial v_3}{\partial v_2}
        \frac{\partial v_2}{\partial x_1}
        =
        \frac{1}{v_2}\frac{\partial v_2}{\partial x_1}.
        \]
        At the given point,
        \[
        \frac{\partial v_3}{\partial x_1}
        =
        \frac{1}{1/2}\cdot \frac{\sqrt{3}}{2}
        =
        \sqrt{3}.
        \]

        Since $v_4 = x_2$, it does not depend on $x_1$, so
        \[
        \frac{\partial v_4}{\partial x_1} = 0.
        \]

        Since $v_5 = e^{v_4}$,
        \[
        \frac{\partial v_5}{\partial x_1}
        =
        \frac{\partial v_5}{\partial v_4}
        \frac{\partial v_4}{\partial x_1}
        =
        e^{v_4}\cdot 0 = 0.
        \]

        Since $v_6 = v_1 v_5$, by the product rule,
        \[
        \frac{\partial v_6}{\partial x_1}
        =
        \frac{\partial v_1}{\partial x_1}v_5 + v_1\frac{\partial v_5}{\partial x_1}.
        \]
        Evaluating,
        \[
        \frac{\partial v_6}{\partial x_1}
        =
        1\cdot e + \frac{\pi}{6}\cdot 0
        =
        e.
        \]

        Since $v_7 = v_3 + v_6$,
        \[
        \frac{\partial v_7}{\partial x_1}
        =
        \frac{\partial v_3}{\partial x_1} + \frac{\partial v_6}{\partial x_1}
        =
        \sqrt{3} + e.
        \]

        Since $f(x_1,x_2)=v_7$, we conclude
        \[
        \frac{\partial f}{\partial x_1}
        =
        \frac{\partial v_7}{\partial x_1}
        =
        \sqrt{3} + e.
        \]
    \end{enumerate}
\end{enumerate}
\end{solution}


\newpage

%%%%%%%%%%%%%%%%%%%%%%%%%%%%%%%%%%%%%%%%%%%%
% Problem 3
%%%%%%%%%%%%%%%%%%%%%%%%%%%%%%%%%%%%%%%%%%%%%

\begin{problem}[Neural Scaling Laws, 50 pts]
In this problem, we look pr an observation that is extremely relevant and widely utilized in industry and research.

Large-scale experiments have revealed a striking observation: the test loss of neural networks follows \textbf{power laws} in model size and dataset size. In this problem, you will discover these laws yourself on a smaller scale, use them to make predictions, and derive the optimal way to allocate a fixed compute budget.
Empirically, for many architectures and datasets, test loss $L$ follows:
\begin{equation}\label{eq:scaling}
    L(N, D) \approx \frac{a}{N^{\alpha}} + \frac{b}{D^{\beta}} + L_{\infty}
\end{equation}
where $N$ is the number of trainable parameters, $D$ is the number of training samples, and $L_\infty$ is the irreducible error (the best loss any model could achieve on this task). On a log-log plot, each term appears as a straight line.
There is no first-principles derivation of \emph{why} this particular functional form holds, it is an empirical observation. However, the form makes intuitive sense through the lens of bias-variance:
\begin{itemize}
    \item $a / N^\alpha$: \textbf{approximation error}. A small model cannot represent the true function well. As we add parameters, this error decreases --- but with diminishing returns (power law, not exponential).
    \item $b / D^\beta$: \textbf{estimation error}. Even a large model cannot learn well from too little data. More data reduces this error, again with diminishing returns.
    \item $L_\infty$: \textbf{irreducible error}. Some loss remains no matter how large the model or dataset --- due to label ambiguity, insufficient input information (e.g., low-resolution images), or inherent noise.
\end{itemize}
\subsection*{Setup: Architecture}
You will use a depth-scaled residual CNN (ResNet) with a fixed base channel width $C = 64$ and a variable number of residual blocks $K$:
\begin{center}
\begin{tabular}{@{}l@{\;}l@{}}
\textbf{Stem:} & \texttt{Conv2d}$(1, C, 3, \text{padding}=1, \text{bias}=\text{False})$ $\to$ \texttt{BatchNorm2d}$(C)$ $\to$ \texttt{ReLU} $\to$ \texttt{MaxPool2d}$(2)$ \\[6pt]
\textbf{$K \times$ ResBlock:} & each block consists of \\[2pt]
& \quad \texttt{Conv2d}$(C, C, 3, \text{padding}=1, \text{bias}=\text{False})$ $\to$ \texttt{BatchNorm2d}$(C)$ $\to$ \texttt{ReLU} \\[2pt]
& \quad $\to$ \texttt{Conv2d}$(C, C, 3, \text{padding}=1, \text{bias}=\text{False})$ $\to$ \texttt{BatchNorm2d}$(C)$ + skip $\to$ \texttt{ReLU} \\[6pt]
\textbf{Head:} & \texttt{AdaptiveAvgPool2d}$(1)$ $\to$ \texttt{Flatten} $\to$ \texttt{Linear}$(C, 10)$
\end{tabular}
\end{center}

Fashion-MNIST images are $28 \times 28$ with 1 channel. After the stem's max-pool the spatial size is $14 \times 14$; the residual blocks preserve this size; the adaptive average pool reduces it to $1 \times 1$. You are free to use \texttt{torch.nn} and \texttt{torchvision} for this problem.
\begin{enumerate}[label=(\alph*)]
%% Part (a): Discover the power laws
\item \textbf{[22 pts]} Train ResNets on Fashion-MNIST at varying scales.
\begin{enumerate}[label=(\roman*)]
    \item \textbf{[10 pts]} Write out the total parameter count $N$ as a function of $K$ (number of residual blocks) for the architecture above. Be careful: each \texttt{Conv2d} with \texttt{bias=False} has $\text{kernel}^2 \times C_{\text{in}} \times C_{\text{out}}$ parameters, each \texttt{BatchNorm2d}$(C)$ has $2C$ parameters (scale and shift), and the linear layer has $\text{in\_features} \times \text{out\_features} + \text{out\_features}$.
    Write the training code for the ResNet. The values for the different $K$ values are provided for your convenience.
    \item \textbf{[6 pts]} Fix $K = 12$ and train on random subsets of size\\ $D \in \{1000, 2000, 5000, 10000, 20000, 30000, 50000\}$. Record the final test loss for each.
    \item \textbf{[6 pts]} Make two log-log plots:
    \begin{itemize}
        \item $\log(L - \hat{L}_\infty)$ vs.\ $\log N$ (at fixed $D$)
        \item $\log(L - \hat{L}_\infty)$ vs.\ $\log D$ (at fixed $N$)
    \end{itemize}
    where $\hat{L}_\infty$ is estimated as the loss of the best configuration \emph{within each sweep} (i.e., the $K=18$ run for the model-size plot, and the $D=50000$ run for the data-size plot). This anchor point is excluded from the fit in each case.
    Fit a line to each plot. Report your estimated exponents $\hat{\alpha}$ and $\hat{\beta}$, the $R^2$ values, as well as your plots, down below.
    Discuss briefly: why do we need to subtract $\hat{L}_\infty$ before taking the log? What would the plot look like if we skipped this step? Feel free to modify the graphing code to see what it would look like, though that is not necessary.
\end{enumerate}
%% Part (b): Optimal compute allocation
\item \textbf{[10 pts]} Suppose training costs $C_{\text{compute}}$ FLOPs, and assume the approximation $C_{\text{compute}} \approx 6ND$ (forward + backward pass $\approx 6$ multiply-adds per parameter per sample).
\begin{enumerate}[label=(\roman*)]
    \item \textbf{[5 pts]} For a fixed compute budget $C_{\text{compute}}$, we want to choose $N$ and $D$ to minimize the loss in Eq.~\eqref{eq:scaling}. Using the constraint $D = C_{\text{compute}} / (6N)$, substitute into the loss and minimize over $N$. Show that the compute-optimal parameter count satisfies:
    $$
    N^*(C_{\text{compute}}) \propto C_{\text{compute}}^{\,\beta / (\alpha + \beta)}
    $$
    and the compute-optimal dataset size satisfies:
    $$
    D^*(C_{\text{compute}}) \propto C_{\text{compute}}^{\,\alpha / (\alpha + \beta)}.
    $$
    \item \textbf{[3 pts]} Plug in your empirically estimated $\hat{\alpha}$ and $\hat{\beta}$ from part (a). If you double your compute budget, by what factor should you increase $N$? By what factor should you increase $D$?
    \item \textbf{[2 pts]} The ``Chinchilla'' result (Hoffmann et al., 2022) found $\alpha \approx \beta$ for large language models, implying $N$ and $D$ should be scaled equally. If instead $\alpha \gg \beta$ (model size matters much more than data), what would the optimal strategy look like? What about $\alpha \ll \beta$?
\end{enumerate}
%% Part (c): Predict, then verify
\item \textbf{[8 pts]} This is the key test of whether scaling laws are actually useful: can you predict the performance of a model you haven't trained yet?
\begin{enumerate}[label=(\roman*)]
    \item \textbf{[4 pts]} Using \emph{only} your fitted scaling law from part (a), predict the test loss at $D = 60000$. Then actually train that model and report the true test loss.
    How close was your prediction? Discuss potential reasons for any discrepancy. \\
    Hint: What problems may arise from using $D=50000$ as your irreducible loss, and what value does it get?
    \item \textbf{[4 pts]} The $K=18$ model trained on the full dataset achieved a test cross-entropy of $0.2195$. Suppose you are limited to the $K=12$ architecture. According to your fitted data scaling law, how much training data would $K=12$ need to match $K=18$'s performance? Is this achievable? What does this tell you about the relationship between model capacity and data?
\end{enumerate}
\end{enumerate}

\end{problem}

\begin{solution}
\begin{enumerate}
    \item[(a)] 
    \begin{enumerate}
        \item We compute the total number of trainable parameters as a function of the number of residual blocks $K$.

        The stem consists of
        \[
        \texttt{Conv2d}(1,64,3,\text{padding}=1,\text{bias}=\text{False}),
        \]
        so the convolution contributes
        \[
        3 \cdot 3 \cdot 1 \cdot 64 = 576
        \]
        parameters. The batch normalization layer contributes
        \[
        2 \cdot 64 = 128
        \]
        parameters. Thus the stem contributes
        \[
        576 + 128 = 704.
        \]

        Each residual block contains two convolutions of the form
        \[
        \texttt{Conv2d}(64,64,3,\text{padding}=1,\text{bias}=\text{False}),
        \]
        and each such convolution contributes
        \[
        3 \cdot 3 \cdot 64 \cdot 64 = 36864
        \]
        parameters. Since there are two of them, the convolutional part contributes
        \[
        2 \cdot 36864 = 73728.
        \]
        Each block also contains two batch normalization layers, each contributing
        \[
        2 \cdot 64 = 128,
        \]
        so the total batch norm contribution per block is
        \[
        2 \cdot 128 = 256.
        \]
        Therefore each residual block contributes
        \[
        73728 + 256 = 73984
        \]
        parameters.

        The head is a linear layer
        \[
        \texttt{Linear}(64,10),
        \]
        so it contributes
        \[
        64 \cdot 10 + 10 = 650
        \]
        parameters.

        Hence the total parameter count is
        \[
        N(K) = 704 + 73984K + 650 = 1354 + 73984K.
        \]

        \item For the fixed architecture $K=12$, the final test losses from the data sweep are:
        \[
        \begin{array}{c|c}
        D & \text{test loss} \\ \hline
        1000 & 0.65339 \\
        2000 & 0.56990 \\
        5000 & 0.44648 \\
        10000 & 0.33728 \\
        20000 & 0.29714 \\
        30000 & 0.28619 \\
        50000 & 0.236778
        \end{array}
        \]
        As expected, the test loss decreases as the amount of training data increases.

        \item We fit the scaling law
        \[
        L(N,D) \approx \frac{a}{N^\alpha} + \frac{b}{D^\beta} + L_\infty.
        \]

        For the model-size sweep, using the $K=18$ run as the anchor $\hat L_\infty$ and excluding that anchor from the fit, we obtain
        \[
        \hat{\alpha} = 1.119719,
        \qquad
        R_N^2 \approx 0.9683.
        \]

        For the data-size sweep, using the $D=50000$ run as the anchor $\hat L_\infty$ and excluding that anchor from the fit, we obtain
        \[
        \hat{\beta} = 0.668627,
        \qquad
        R_D^2 \approx 0.9763.
        \]

        We need to subtract $\hat L_\infty$ before taking logs because the empirical law is a power law for the \emph{gap}
        \[
        L - \hat L_\infty,
        \]
        not for $L$ itself. In other words, the approximately linear relationship on log-log axes is
        \[
        \log(L-\hat L_\infty) \approx \log a - \alpha \log N
        \]
        or
        \[
        \log(L-\hat L_\infty) \approx \log b - \beta \log D.
        \]
        If we skipped the subtraction, then as the loss approached its floor the log-log plot would bend and flatten rather than appearing approximately linear.
        \begin{figure}
            \centering
            \includegraphics[width=0.8\linewidth]{output3aiii.pdf}
            \label{fig:placeholder}
        \end{figure}
    \end{enumerate}

    \item[(b)]
    \begin{enumerate}
        \item Suppose the compute budget satisfies
        \[
        C_{\text{compute}} \approx 6ND.
        \]
        Then
        \[
        D = \frac{C_{\text{compute}}}{6N}.
        \]
        Substituting this into the scaling law gives
        \[
        L(N) \approx \frac{a}{N^\alpha} + \frac{b}{\left(C_{\text{compute}}/(6N)\right)^\beta} + L_\infty
        = aN^{-\alpha} + b\left(\frac{6N}{C_{\text{compute}}}\right)^\beta + L_\infty.
        \]

        Differentiate with respect to $N$:
        \[
        \frac{dL}{dN}
        =
        -\alpha a N^{-\alpha-1}
        +
        \beta b\,6^\beta C_{\text{compute}}^{-\beta} N^{\beta-1}.
        \]
        Setting this equal to zero,
        \[
        \alpha a N^{-\alpha-1}
        =
        \beta b\,6^\beta C_{\text{compute}}^{-\beta} N^{\beta-1}.
        \]
        Rearranging,
        \[
        N^{\alpha+\beta}
        \propto
        C_{\text{compute}}^\beta.
        \]
        Therefore,
        \[
        N^*(C_{\text{compute}})
        \propto
        C_{\text{compute}}^{\beta/(\alpha+\beta)}.
        \]

        Since
        \[
        D = \frac{C_{\text{compute}}}{6N},
        \]
        we also get
        \[
        D^*(C_{\text{compute}})
        \propto
        C_{\text{compute}}^{1-\beta/(\alpha+\beta)}
        =
        C_{\text{compute}}^{\alpha/(\alpha+\beta)}.
        \]

        \item Using
        \[
        \hat\alpha = 1.119719,
        \qquad
        \hat\beta = 0.668627,
        \]
        we have
        \[
        \hat\alpha + \hat\beta = 1.788346.
        \]
        If compute doubles, then
        \[
        \frac{N_{\text{new}}}{N_{\text{old}}}
        =
        2^{\hat\beta/(\hat\alpha+\hat\beta)}
        =
        2^{0.668627/1.788346}
        \approx 1.296,
        \]
        and
        \[
        \frac{D_{\text{new}}}{D_{\text{old}}}
        =
        2^{\hat\alpha/(\hat\alpha+\hat\beta)}
        =
        2^{1.119719/1.788346}
        \approx 1.543.
        \]

        So if we double the compute budget, we should increase
        \[
        N \text{ by about } 1.30\times
        \qquad\text{and}\qquad
        D \text{ by about } 1.54\times.
        \]

        \item If $\alpha \gg \beta$, then model size matters much more than data. In that case, the compute-optimal strategy would place more of the extra compute into increasing $N$ rather than $D$.

        If $\alpha \ll \beta$, then data matters much more than model size. In that case, the compute-optimal strategy would place more of the extra compute into increasing $D$ rather than $N$.
    \end{enumerate}

    \item[(c)]
    \begin{enumerate}
        \item Using only the fitted data scaling law from part (a), the predicted test loss at
        \[
        D=60000
        \]
        is
        \[
        \hat L(60000) = 0.268425.
        \]
        After actually training the $K=12$ model on the full dataset, the true test loss is
        \[
        L_{\text{true}}(60000)=0.257493.
        \]
        Therefore the absolute prediction error is
        \[
        |0.268425 - 0.257493| = 0.010932.
        \]

        This prediction is fairly close. One possible source of discrepancy is that in the fit we used the $D=50000$ result as an estimate of the irreducible loss $\hat L_\infty$, but this is not actually a true irreducible loss. It is only the best observed loss in the sweep, so it may still be above the true floor. In addition, optimization noise, finite-sample noise, random subset selection, and the fact that the power-law relationship is only approximate can all contribute to prediction error.

        \item The $K=18$ model trained on the full dataset achieved a test cross-entropy of
        \[
        0.2195.
        \]
        For the $K=12$ architecture, the fitted data scaling law has the form
        \[
        L(D) \approx \hat L_\infty + \hat b D^{-\hat\beta}.
        \]
        Here, from the data sweep we used
        \[
        \hat L_\infty = 0.236778.
        \]
        To match the $K=18$ performance, we would need
        \[
        0.2195 = \hat L_\infty + \hat b D^{-\hat\beta}.
        \]
        But this is impossible, because
        \[
        0.2195 < 0.236778 = \hat L_\infty.
        \]
        Since the target loss lies below the fitted asymptotic floor for the $K=12$ model, there is no finite positive value of $D$ that solves this equation under the fitted scaling law.

        Therefore, according to the fitted law, the $K=12$ architecture cannot match the $K=18$ model's performance by adding more data alone. This suggests that model capacity and data are not interchangeable indefinitely: once a smaller model hits its apparent performance floor, additional data cannot fully compensate for insufficient capacity.
    \end{enumerate}
\end{enumerate}
\end{solution}

\newpage


%%%%%%%%%%%%%%%%%%%%%%%%%%%%%%%%%%%%%%%%%%%%%
% Name and Calibration
%%%%%%%%%%%%%%%%%%%%%%%%%%%%%%%%%%%%%%%%%%%%%

\textbf{Name}: Ibrahim Khaliliya

\textbf{Collaborators and Resources}: Claude was used for LaTeX styling issues, especially with the figures.

\end{document}
